% ==========================================================================
%  Chapter 45 — Phase 0 Implementation: PETSc RAII Wrappers
% ==========================================================================

\chapter{Phase 0: PETSc RAII and Error Discipline}
\label{ch:phase0-petsc-raii}

This chapter documents the implementation of Phase~0 from the master plan
(Chapter~\ref{ch:convergence-master-plan}): wrapping every PETSc handle in
a move-only RAII class and replacing manual \texttt{Destroy} calls
throughout the codebase.

%% ------------------------------------------------------------------
\section{Motivation}
\label{sec:phase0-motivation}

PETSc objects (\texttt{Vec}, \texttt{Mat}, \texttt{KSP}, \texttt{SNES},
\texttt{TS}, \texttt{PetscSection}) follow a C-style
\texttt{Create}/\texttt{Destroy} lifecycle.  The existing code managed
these with explicit destructor logic:

\begin{verbatim}
~Model() {
    if (Kt_)           MatDestroy(&Kt_);
    if (nodal_forces_) VecDestroy(&nodal_forces_);
    // ... 10+ more lines
}
\end{verbatim}

This approach has three problems:

\begin{enumerate}
  \item \textbf{Leak on early return.}  If a method calls
        \texttt{VecCreate} then exits before reaching the destructor
        (thrown exception, early \texttt{return}, failed assertion),
        the handle is leaked.
  \item \textbf{Error in cleanup.}  Forgetting to add a new member to the
        destructor is an easy mistake as classes grow.
  \item \textbf{Unchecked errors.}  PETSc routines return
        \texttt{PetscErrorCode}, but the codebase silently discarded it.
        A failed allocation or invalid argument went unnoticed until
        a mysterious crash later.
\end{enumerate}

%% ------------------------------------------------------------------
\section{Design: Thin Wrappers with Implicit Conversion}
\label{sec:phase0-design}

Each wrapper follows the same pattern.  We illustrate with \texttt{OwnedVec};
the others (\texttt{OwnedMat}, \texttt{OwnedKSP}, etc.) are identical up to
the handle type and the corresponding \texttt{Destroy} function.

\subsection{Invariants}

\begin{itemize}
  \item \textbf{Move-only.}  Copy construction and copy assignment are
        deleted.  Move transfers ownership and nullifies the source.
  \item \textbf{Implicit conversion.}  \texttt{operator Vec() const}
        returns the raw handle, so the wrapper can be passed directly to
        any PETSc function that takes a \texttt{Vec} by value.
  \item \textbf{Pointer-to-member.}  \texttt{ptr()} returns
        \texttt{Vec*}, for use with PETSc creation functions that take
        an output pointer (\texttt{VecCreate(..., \&v)}).
  \item \textbf{Boolean test.}  \texttt{explicit operator bool()} returns
        \texttt{true} iff the handle is non-null (guards conditional creation).
  \item \textbf{Destructor calls \texttt{Destroy}.}  If the handle is
        non-null, the destructor calls the corresponding PETSc destroy
        function.  This is the entire point of RAII.
\end{itemize}

\subsection{Interface (abridged)}

\begin{verbatim}
namespace petsc {

class OwnedVec {
    Vec v_{nullptr};
public:
    OwnedVec() = default;
    explicit OwnedVec(Vec v) : v_{v} {}

    ~OwnedVec() { if (v_) VecDestroy(&v_); }

    OwnedVec(OwnedVec&& o) noexcept : v_{o.v_} { o.v_ = nullptr; }
    OwnedVec& operator=(OwnedVec&& o) noexcept {
        if (this != &o) { reset(); v_ = o.v_; o.v_ = nullptr; }
        return *this;
    }

    OwnedVec(const OwnedVec&)            = delete;
    OwnedVec& operator=(const OwnedVec&) = delete;

    operator Vec() const { return v_; }
    Vec  get()     const { return v_; }
    Vec* ptr()           { reset(); return &v_; }

    explicit operator bool() const { return v_ != nullptr; }

    Vec release() { Vec t = v_; v_ = nullptr; return t; }

    void reset() {
        if (v_) { VecDestroy(&v_); v_ = nullptr; }
    }
};

} // namespace petsc
\end{verbatim}

\paragraph{Key design choice: \texttt{ptr()} resets before returning.}
This ensures that calling \texttt{DMCreateGlobalVector(dm, v.ptr())}
does \emph{not} leak a previously held handle.  The pattern mirrors
\texttt{std::unique\_ptr::reset()} before \texttt{release()}.

\subsection{Full Wrapper Set}

\begin{center}
\begin{tabular}{lll}
\toprule
\textbf{Wrapper} & \textbf{PETSc type} & \textbf{Destroy function} \\
\midrule
\texttt{OwnedVec}     & \texttt{Vec}           & \texttt{VecDestroy} \\
\texttt{OwnedMat}     & \texttt{Mat}           & \texttt{MatDestroy} \\
\texttt{OwnedKSP}     & \texttt{KSP}           & \texttt{KSPDestroy} \\
\texttt{OwnedSNES}    & \texttt{SNES}          & \texttt{SNESDestroy} \\
\texttt{OwnedTS}      & \texttt{TS}            & \texttt{TSDestroy} \\
\texttt{OwnedSection} & \texttt{PetscSection}  & \texttt{PetscSectionDestroy} \\
\bottomrule
\end{tabular}
\end{center}

All are located in \texttt{src/petsc/} and collected by the convenience
header \texttt{src/petsc/PetscRaii.hh}.

%% ------------------------------------------------------------------
\section{Error Checking Macro}
\label{sec:phase0-check}

A single macro wraps every PETSc call:

\begin{verbatim}
#define FALL_N_PETSC_CHECK(expr)                                   \
    do {                                                            \
        PetscErrorCode ierr_ = (expr);                              \
        if (ierr_) throw petsc::PetscError(__FILE__, __LINE__,      \
                                           #expr, ierr_);          \
    } while (0)
\end{verbatim}

The exception type \texttt{petsc::PetscError} stores file, line, the
stringified expression, and the error code.  In release mode the macro
compiles to a no-op unless \texttt{FALL\_N\_PETSC\_CHECK\_RELEASE} is defined.

Defined in \texttt{src/petsc/check.hh}.

%% ------------------------------------------------------------------
\section{Refactored Classes}
\label{sec:phase0-refactored}

\subsection{Model}

Before:
\begin{verbatim}
Mat Kt_{nullptr};
Vec nodal_forces_{nullptr};
Vec current_state_{nullptr};
Vec imposed_solution_{nullptr};
PetscSection dof_section{nullptr};

~Model() {
    if (Kt_)              MatDestroy(&Kt_);
    if (nodal_forces_)    VecDestroy(&nodal_forces_);
    if (current_state_)   VecDestroy(&current_state_);
    if (imposed_solution_) VecDestroy(&imposed_solution_);
    /* ... 4 more lines ... */
}
\end{verbatim}

After:
\begin{verbatim}
petsc::OwnedMat     Kt_{};
petsc::OwnedVec     nodal_forces_{};
petsc::OwnedVec     current_state_{};
petsc::OwnedVec     imposed_solution_{};
petsc::OwnedSection dof_section_{};

~Model() = default;
\end{verbatim}

All creation calls were changed from \texttt{DMCreateLocalVector(dm, \&v)}
to \texttt{FALL\_N\_PETSC\_CHECK(DMCreateLocalVector(dm, v.ptr()))}.

Public accessors return the raw handle via implicit conversion:
\texttt{return current\_state\_.get();} (or simply \texttt{return current\_state\_;}).

\subsection{LinearAnalysis}

Raw \texttt{KSP}, \texttt{Mat}, and \texttt{Vec} members replaced with
\texttt{OwnedKSP}, \texttt{OwnedMat}, \texttt{OwnedVec}.  Manual destructor
removed.

\subsection{NonlinearAnalysis}

Raw \texttt{SNES}, \texttt{Mat}, and \texttt{Vec} members replaced.
Local temporaries (\texttt{f\_full}, \texttt{U\_snap}) in the bisection
algorithm are now \texttt{OwnedVec} instances whose lifetime is
scope-bounded---no explicit \texttt{VecDestroy} needed.

\subsection{DynamicAnalysis}

Raw \texttt{TS}, \texttt{Mat} (M\_, C\_, J\_, K0\_), and \texttt{Vec}
(U\_, V\_, F\_work\_, f\_ext\_work\_) replaced.  The
\texttt{GroundMotionInfo} struct's \texttt{influence} field changed from
raw \texttt{Vec} to \texttt{petsc::OwnedVec}.  The 15-line manual
destructor collapses to \texttt{= default}.

\paragraph{Static callbacks.}  The PETSc TS callbacks
(\texttt{I2Function}, \texttt{I2Jacobian}, \texttt{Monitor}) access
owned members via \texttt{self->M\_}, which works transparently thanks to
the implicit \texttt{operator Mat()} conversion---the PETSc C API
receives the raw handle by value.

%% ------------------------------------------------------------------
\section{File Layout}

\dirtree{%
  .1 src/.
  .2 petsc/.
  .3 check.hh\DTcomment{FALL\_N\_PETSC\_CHECK macro + PetscError}.
  .3 OwnedVec.hh.
  .3 OwnedMat.hh.
  .3 OwnedKSP.hh.
  .3 OwnedSNES.hh.
  .3 OwnedTS.hh.
  .3 OwnedSection.hh.
  .3 PetscRaii.hh\DTcomment{convenience include-all}.
  .2 model/.
  .3 Model.hh\DTcomment{refactored}.
  .2 analysis/.
  .3 LinearAnalysis.hh\DTcomment{refactored}.
  .3 NLAnalysis.hh\DTcomment{refactored}.
  .3 DynamicAnalysis.hh\DTcomment{refactored}.
}

%% ------------------------------------------------------------------
\section{Validation Criteria}

Phase~0 is considered complete when:

\begin{enumerate}
  \item The project compiles with no warnings related to PETSc handles.
  \item No class contains a manual \texttt{Destroy} call for owned
        PETSc objects.
  \item Every PETSc creation/query call that can fail is wrapped in
        \texttt{FALL\_N\_PETSC\_CHECK}.
  \item Existing tests (\texttt{fall\_n\_tests},
        \texttt{fall\_n\_continuum\_test}, \texttt{fall\_n\_dynamic\_test},
        etc.) pass without regressions.
  \item LeakSanitizer reports no new PETSc-related leaks
        (existing suppressions in \texttt{lsan.supp} remain valid).
\end{enumerate}
